\documentclass[11pt,a4paper]{article}
\usepackage[margin=25mm]{geometry}
\usepackage{iftex}
\ifPDFTeX
  \usepackage[T1]{fontenc}
  \usepackage[utf8]{inputenc}
\else
  \usepackage{fontspec}
  \setmainfont{cmunrm.otf}[
    BoldFont=cmunbx.otf,
    ItalicFont=cmunti.otf,
    BoldItalicFont=cmunbi.otf,
    Ligatures=TeX]
\fi
\usepackage[english]{babel}
\usepackage{amsmath,amssymb,amsthm}
\usepackage{mathrsfs}
\usepackage{microtype}
\usepackage{enumitem}
\usepackage[unicode]{hyperref}
\hypersetup{
  pdftitle={Unbounded Orbit-Series and Reductivity Indices of Finite Quandles},
  pdfauthor={Danil Pokulevskii},
  pdflang={en-US},
  colorlinks=true,linkcolor=blue,citecolor=blue,urlcolor=blue}
\setlength{\emergencystretch}{2em}
\setlist[enumerate]{itemsep=2pt,topsep=5pt}
\frenchspacing
\newtheorem{theorem}{Theorem}[section]
\newtheorem{lemma}[theorem]{Lemma}
\newtheorem{proposition}[theorem]{Proposition}
\newtheorem{corollary}[theorem]{Corollary}
\theoremstyle{definition}
\newtheorem{definition}[theorem]{Definition}
\theoremstyle{remark}
\newtheorem{remark}[theorem]{Remark}
\renewcommand{\proofname}{Proof}
\newcommand{\R}{\mathcal R}
\newcommand{\LR}{\mathcal{LR}}
\newcommand{\tOS}{t\mathcal{OS}}
\newcommand{\Inn}{\operatorname{Inn}}
\newcommand{\Orb}{\operatorname{Orb}}
\newcommand{\cl}{\operatorname{class}}
\newcommand{\ad}{\operatorname{ad}}
\newcommand{\Z}{\mathbb Z}
\newcommand{\F}{\mathbb F}
\newcommand{\conj}{\triangleright}

\title{Unbounded Orbit-Series and Reductivity Indices\\
of Finite Quandles}
\author{Danil Pokulevskii\\[3pt]
  \small Novosibirsk State Technical University\\
  \small Novosibirsk, Russia\\
  \small\href{mailto:555tery@gmail.com}{\texttt{555tery@gmail.com}}}
\date{October 7, 2026\\[2pt]\small Working preprint}

\begin{document}
\maketitle
\begin{abstract}
For a finite quandle, let \(k\), \(j\), and \(i\) denote the
least positive indices of local reductivity, trivialization of
orbit series, and reductivity, respectively.
We show that \(j\) cannot be bounded by a function of \(k\),
and \(i\) cannot be bounded by a function of \(j\).
For every \(n\ge3\), we construct a finite quandle \(Q_n\) with
\(k(Q_n)=2\) and \(j(Q_n)\ge\lfloor\log_3n\rfloor+1\).
For every \(m\ge2\), we construct a finite quandle \(P_m\) with
\((k(P_m),j(P_m),i(P_m))=(2,2,m)\).
The first family is obtained from the group of Noce--Tracey--Traustason
using a filtration by label size; the second is obtained from finite
quotients of the reduced free group that preserve its exact
nilpotency class. All passages from groups to quandles,
including the shift between the nilpotency class of the inner
automorphism group and the reductivity index, are proved explicitly.
This gives negative answers to both parts of Question~5.25
of Bonatto--Crans--Nasybullov--Whitney.
\end{abstract}

\noindent\textbf{Keywords:} finite quandle, local reductivity,
orbit series, reduced free group, Engel element.

\section{Introduction and main results}

The classes \(\LR_n\) and \(\tOS_n\) were introduced in~\cite{BCNW},
where their relationship with the classes \(\R_n\) was investigated.
For finite quandles, these finite-depth conditions are equivalent
at the level of the unions of the respective classes
\cite[Corollary~5.19]{BCNW}.
Question~5.25 asks whether orbit-series depth can be bounded
in terms of the local reductivity index, and whether the reductivity
index can be bounded in terms of orbit-series depth.
The numbering of the results in~\cite{BCNW} used here follows
arXiv version v3, dated August 18, 2020.

For a quandle whose corresponding indices exist, set
\[
 \begin{aligned}[c]
 k(Q)&=\min\{n\ge1:Q\in\LR_n\},\\
 j(Q)&=\min\{n\ge1:Q\in\tOS_n\},\\
 i(Q)&=\min\{n\ge1:Q\in\R_n\}.
 \end{aligned}
\]
The classes are defined in Section~\ref{sec:definitions}.

\begin{theorem}\label{thm:first}
For every integer \(n\ge3\), there is a finite quandle \(Q_n\)
such that
\[
 k(Q_n)=2,\qquad j(Q_n)\ge\lfloor\log_3n\rfloor+1.
\]
\end{theorem}

\begin{theorem}\label{thm:second}
For every integer \(m\ge2\), there is a finite quandle \(P_m\)
such that
\[
 (k(P_m),j(P_m),i(P_m))=(2,2,m).
\]
The quandle \(P_m\) can be chosen to have exactly \(m\) inner orbits,
each of which is trivial as a quandle.
\end{theorem}

\begin{corollary}\label{cor:nobounds}
There is no function \(f:\mathbb N_{>0}\to\mathbb N_{>0}\)
such that \(j(Q)\le f(k(Q))\) for every finite quandle
with defined indices. Nor is there a function \(g\)
such that \(i(Q)\le g(j(Q))\) for all such quandles.
\end{corollary}
\begin{proof}
The first family has \(k=2\) and unbounded \(j\).
The second has \(j=2\) and unbounded \(i\).
\end{proof}

Theorems~\ref{thm:first} and~\ref{thm:second} are proved in
Sections~\ref{sec:first} and~\ref{sec:second}.
We use two known group constructions:
the group of Noce--Tracey--Traustason~\cite{NTT}
and the reduced free group as treated by Darn\'e~\cite{Darne}.
The group-theoretic input results are stated separately from
the consequences for quandles established here.

\section{Definitions and a reductivity criterion}\label{sec:definitions}

We use the left convention for quandles~\cite{BCNW}:
a quandle is a nonempty set \(Q\) with an operation \(\conj\) satisfying
\[
 a\conj a=a,\qquad L_a:b\mapsto a\conj b\ \text{is bijective},
 \qquad a\conj(b\conj c)=(a\conj b)\conj(a\conj c).
\]
The group \(\Inn(Q)=\langle L_a:a\in Q\rangle\)
is its inner automorphism group.
Subquandles are also closed under inverse left translations.
The trivial operation is given by \(a\conj b=b\).

Write
\[
 T_n(a;c_1,\ldots,c_n)
 =((a\conj c_1)\conj\cdots)\conj c_n
\]
for the left-associated expression with \(n\) operations.
The class \(\R_n\) is defined by the identity
\[
 T_n(a;c_1,\ldots,c_n)=T_n(b;c_1,\ldots,c_n)
\]
for all arguments; the class \(\LR_n\) is defined by
\[
 T_n(a;\underbrace{b,\ldots,b}_{n})=b.
\]

An orbit series is defined recursively by
\[
 Q_0=Q,\qquad Q_{r+1}=\Orb(x_r,Q_r),\qquad x_r\in Q_r.
\]
At each step, the inner automorphism group of the current subquandle
is used. The class \(\tOS_n\) consists of quandles for which every
such series reaches a singleton quandle no later than step \(n\).
Thus \(j\) measures the time to reach a singleton,
rather than merely the time to stabilization.
All three classes with index \(1\) coincide with the class of trivial quandles.
Moreover, \(Q\in\tOS_2\) if and only if each inner orbit of \(Q\)
is trivial as a quandle; such an orbit may contain many elements.

For groups, we use the conventions
\[
 [a,b]=a^{-1}b^{-1}ab,\qquad c^g=g^{-1}cg,\qquad
 \gamma_1(H)=H,\quad \gamma_{r+1}(H)=[\gamma_r(H),H].
\]
The trivial group has class zero.
Otherwise, exact class \(c\) means
\(\gamma_{c+1}(H)=1\), \(\gamma_c(H)\ne1\).

\begin{lemma}\label{lem:orbits}
Let \(S\) be a nonempty subset of a group, closed under
conjugation by its own elements and their inverses.
Then \(a\conj b=a^{-1}ba\) makes \(S\) a quandle, and for \(a\in S\)
\[
 \Orb(a,S)=a^{\langle S\rangle}.
\]
\end{lemma}
\begin{proof}
Idempotency and bijectivity are immediate, and
\[
 (a^{-1}ba)^{-1}(a^{-1}ca)(a^{-1}ba)=a^{-1}b^{-1}cba
\]
proves self-distributivity.
Left translations and their inverses are conjugations by
elements of \(S\) and their inverses. They generate the action
of \(\langle S\rangle\), giving the orbit formula.
\end{proof}

\begin{proposition}\label{prop:criterion}
For every nonempty quandle \(Q\) and \(n\ge1\),
\[
 Q\in\R_n\quad\Longleftrightarrow\quad
 \gamma_n(\Inn(Q))=1.
\]
In particular, if \(\Inn(Q)\) has exact finite class \(c\),
then \(i(Q)=c+1\).
\end{proposition}
\begin{proof}
Set \(\mathcal L(Q)=\{L_a:a\in Q\}\).
Self-distributivity gives
\[
 L_{a\conj b}=L_aL_bL_a^{-1}.
\]
Thus \(a\mapsto L_a\) is an epimorphism of quandles
when \(\mathcal L(Q)\) is equipped with \(s\conj t=sts^{-1}\).
This direction of conjugation is opposite to the convention
in Lemma~\ref{lem:orbits}.

For \(n\ge1\), the identity \(\R_{n+1}\) for \(Q\)
is equivalent to the identity \(\R_n\) for \(\mathcal L(Q)\):
equality of the two expressions after the final operation
for every final argument means equality of their left translations.
Also, \(Q\in\R_1\) if and only if \(\mathcal L(Q)\) is a singleton.
Indeed, if all \(L_a\) equal \(f\), then idempotency gives
\(f(a)=a\) for every \(a\).
For the iterates \(\mathcal L_0(Q)=Q\),
\(\mathcal L_{r+1}(Q)=\mathcal L(\mathcal L_r(Q))\), we obtain
\[
 Q\in\R_n\quad\Longleftrightarrow\quad
 |\mathcal L_n(Q)|=1.
\]

Let \(H=\Inn(Q)\). The set \(\mathcal L(Q)\) generates \(H\)
and is preserved by conjugations in \(H\), since \(hL_ah^{-1}=L_{h(a)}\).
The action of \(H\) on this set has image
\(\Inn(\mathcal L(Q))\) and kernel \(C_H(\mathcal L(Q))=Z(H)\).
Consequently,
\[
 \Inn(\mathcal L(Q))\cong H/Z(H).
\]
For the upper central series \(Z_0(H)=1\),
\(Z_{r+1}(H)/Z_r(H)=Z(H/Z_r(H))\), induction gives
\[
 \Inn(\mathcal L_r(Q))\cong H/Z_r(H).
\]
The quandle \(\mathcal L_{n-1}(Q)\) is trivial if and only if
\(\mathcal L_n(Q)\) is a singleton; hence the latter condition
is equivalent to \(H=Z_{n-1}(H)\).

Finally, \(H=Z_c(H)\) is equivalent to \(\gamma_{c+1}(H)=1\).
This follows by induction on \(c\), starting with \(c=0\):
the condition \(\gamma_{c+1}(H)=1\) is equivalent to
\(\gamma_c(H)\le Z(H)\), or \(\gamma_c(H/Z(H))=1\).
Applying induction to \(H/Z(H)\) gives \(H=Z_c(H)\).
Taking \(c=n-1\) proves the criterion, including \(n=1\).
\end{proof}

\begin{remark}
The criterion corresponds to \cite[Proposition~3.2]{BCNW},
but is expressed by vanishing of a term of the series to distinguish
a bound on the class from its exact value.
\end{remark}

\section{First family: unbounded depth}\label{sec:first}

\subsection{Group-theoretic input}

We use the group \(\mathscr G\) over \(\F_2\) constructed in~\cite{NTT}.
Write its generators as
\[
 X=1+\ad(x),\quad U_A=1+\ad(u_A),\quad
 V_A=1+\ad(v_A),\quad W_A=1+\ad(w_A),
\]
where \(A\) is a nonempty finite subset of \(\mathbb N\).
The operators and their action are defined in~\cite[Section~2]{NTT}.
All these generators are involutions.
Each of the three families \(U,V,W\) generates an elementary abelian group.
If \(A\cap B=\varnothing\), put \(A\sqcup B=A\cup B\);
otherwise put \(A\sqcup B=\varnothing\).
The empty label corresponds to the identity.
By \cite[Lemma~2.3]{NTT},
\begin{equation}\label{eq:nttcomm}
 \begin{aligned}[c]
 [U_A,V_B]&=U_{A\sqcup B},&
 [V_A,W_B]&=W_{A\sqcup B},\\
 [W_A,U_B]&=V_{A\sqcup B},&
 [W_A,X]&=U_A,
 \end{aligned}
\end{equation}
and \(X\) commutes with all \(U_A,V_A\).
By \cite[Proposition~2.4]{NTT}, every element has a unique
normal form \(X^\epsilon uvw\), where \(\epsilon\in\{0,1\}\)
and \(u,v,w\) belong to the corresponding elementary abelian groups.
Finally, the proof of \cite[Theorem~2.5]{NTT} establishes
\begin{equation}\label{eq:pairengel}
 [[X^g,X],X]=1\qquad(g\in\mathscr G).
\end{equation}

Every \(U_A\) is nontrivial.
To see this in the original operator construction,
choose a singleton \(B\) disjoint from \(A\):
the operator \(\ad(u_A)\) sends \(w_B\) to the nonzero \(v_{A\cup B}\).

For \(n\ge1\), set
\[
 G_n=\langle X,U_A,V_A,W_A:\varnothing\ne A\subseteq[n]\rangle
 \le\mathscr G,\qquad [n]=\{1,\ldots,n\}.
\]
When a word is collected into normal form, new labels arise
only through \(\sqcup\), and therefore remain subsets of \([n]\).
It follows that \(G_n\) is finite and
\begin{equation}\label{eq:sizefirst}
 |G_n|\le2^{\,1+3(2^n-1)}.
\end{equation}
Here the labels themselves are restricted, rather than only their sizes.

\subsection{The filtration and normal closures}

For an integer \(s\ge1\), define the subgroup
\[
 F_s=\langle X,\ U_A\ (|A|\ge s),\
 V_A\ (|A|\ge2s),\ W_A\ (|A|\ge3s)\rangle\le G_n,
\]
using only labels in \([n]\).
Its elements are precisely the normal forms with
\[
 u\in\langle U_A:|A|\ge s\rangle,\quad
 v\in\langle V_A:|A|\ge2s\rangle,\quad
 w\in\langle W_A:|A|\ge3s\rangle.
\]
Indeed, the only correction terms in the collection process satisfy
\[
 \begin{aligned}[c]
 [X,W_{\ge3s}]&\subseteq U_{\ge3s},&
 [U_{\ge s},V_{\ge2s}]&\subseteq U_{\ge3s},\\
 [W_{\ge3s},U_{\ge s}]&\subseteq V_{\ge4s},&
 [V_{\ge2s},W_{\ge3s}]&\subseteq W_{\ge5s}.
 \end{aligned}
\]
The threshold does not decrease: disjoint labels give the sum
of their sizes, whereas intersecting labels give the identity.
Hereafter, \(U_{\ge t}\) and the analogous symbols denote subgroups
generated by labels of the corresponding size.

\begin{lemma}\label{lem:filtration}
For every \(s\ge1\),
\[
 \langle X^{F_s}\rangle=F_{3s}.
\]
Moreover, \(\langle X^{G_n}\rangle=F_1\).
\end{lemma}
\begin{proof}
First, we check that \(F_{3s}\trianglelefteq F_s\).
In the table, rows correspond to generators of \(F_{3s}\)
and columns to generators of \(F_s\); an entry indicates the layer
containing a possible nontrivial commutator:
\[
 \begin{array}{c|cccc}
 &X&U_{\ge s}&V_{\ge2s}&W_{\ge3s}\\ \hline
 X&1&1&1&U_{\ge3s}\\
 U_{\ge3s}&1&1&U_{\ge5s}&V_{\ge6s}\\
 V_{\ge6s}&1&U_{\ge7s}&1&W_{\ge9s}\\
 W_{\ge9s}&U_{\ge9s}&V_{\ge10s}&W_{\ge11s}&1
 \end{array}
\]
All entries lie in \(F_{3s}\).
Since the generators are involutions, this check ensures
invariance under conjugation by them in both directions.
The subgroup contains \(X\), so
\(\langle X^{F_s}\rangle\le F_{3s}\).

For the reverse inclusion, set \(K=\langle X^{F_s}\rangle\).
It is normal in \(F_s\) and contains \(X\).
If \(|A|\ge3s\), then \(W_A\in F_s\), so
\([W_A,X]=U_A\in K\).
Every label \(T\) of size at least \(6s\) can be partitioned
as \(A\sqcup B\) with \(|A|,|B|\ge3s\);
then \([W_A,U_B]=V_T\in K\).
For \(|T|\ge9s\), take a partition with
\(|A|\ge6s,\ |B|\ge3s\).
The previous step gives \(V_A\in K\), and hence
\([V_A,W_B]=W_T\in K\).
These are all the generators of \(F_{3s}\).

For the initial step, we check that \(F_1\trianglelefteq G_n\).
Possible nontrivial commutators between generators lie in the layers
\[
 \begin{aligned}[c]
 [X,W_{\ge1}]&\subseteq U_{\ge1},&
 [U_{\ge1},V_{\ge1}]&\subseteq U_{\ge2},\\
 [U_{\ge1},W_{\ge1}]&\subseteq V_{\ge2},&
 [V_{\ge2},W_{\ge1}]&\subseteq W_{\ge3},\\
 [W_{\ge3},U_{\ge1}]&\subseteq V_{\ge4},&
 [W_{\ge3},V_{\ge1}]&\subseteq W_{\ge4},\\
 [W_{\ge3},X]&\subseteq U_{\ge3}.
 \end{aligned}
\]
They all lie in \(F_1\); the others are the identity
or are obtained by inverting a commutator.
Thus \(\langle X^{G_n}\rangle\le F_1\).
Conversely, \([W_A,X]=U_A\) puts all \(U_A\) in the normal closure;
partitions into two nonempty parts and \([W_A,U_B]\)
give all \(V_T\) with \(|T|\ge2\);
partitions with \(|A|\ge2,\ |B|\ge1\) and \([V_A,W_B]\)
give all \(W_T\) with \(|T|\ge3\).
\end{proof}

\subsection{The quandle and its depth}

\begin{proof}[Proof of Theorem~\ref{thm:first}]
Let \(C_n=X^{G_n}\), and equip \(Q_n=C_n\)
with the operation \(a\conj b=a^{-1}ba\).
This is a finite quandle by Lemma~\ref{lem:orbits}.
For \(a=X^g,\ b=X^h\), conjugation by \(h^{-1}\)
sends the pair to \(X^{gh^{-1}},X\).
Thus \eqref{eq:pairengel} gives
\([[a,b],b]=1\) for all \(a,b\in C_n\).
Interchanging \(a,b\) gives \([[b,a],a]=1\).
Since \([b,a]=[a,b]^{-1}\), the commutator \([a,b]\)
centralizes both generators. The subgroup \(\langle a,b\rangle\)
has class at most two.
In particular, \(a^{-1}ba\) commutes with \(b\), so
\[
 (a\conj b)\conj b=b.
\]
Hence \(Q_n\in\LR_2\).

For \(n\ge3\), choose \(|A|=3\).
Then \(W_A\in F_1=\langle C_n\rangle\)
and \([W_A,X]=U_A\ne1\).
By Lemma~\ref{lem:orbits}, the orbit of \(X\) in \(Q_n\)
is nontrivial, so \(Q_n\) is not trivial and \(k(Q_n)=2\).
The index \(j(Q_n)\) exists:
by \cite[Corollary~5.19]{BCNW}, a finite quandle in \(\LR_2\)
belongs to \(\tOS\), meaning that every orbit series
reaches a singleton quandle.
The length of such a series in a finite quandle is bounded:
all inclusions before the singleton term are strict,
since equality of consecutive terms would mean that
the current term is connected and cannot decrease further.
Therefore \(Q_n\in\bigcup_{t\ge1}\tOS_t\).

Choosing the anchor \(X\) at every step, set
\[
 S_0=C_n,\qquad N_r=\langle S_r\rangle,\qquad
 S_{r+1}=X^{N_r}.
\]
This is an orbit series by Lemma~\ref{lem:orbits}.
Lemma~\ref{lem:filtration} gives
\[
 N_r=F_{3^r}\quad(r\ge0),\qquad
 S_r=X^{F_{3^{r-1}}}\quad(r\ge1).
\]
If \(r\ge1\) and \(3^r\le n\), choose \(|A|=3^r\).
Then \(W_A\in F_{3^{r-1}}\), so \(X,X^{W_A}\in S_r\);
these elements are distinct because \([W_A,X]=U_A\ne1\).
Consequently, \(S_r\) is not a singleton.
For \(r=\lfloor\log_3n\rfloor\), the series does not reach
a singleton at any of the steps \(0,\ldots,r\),
and \(j(Q_n)\ge r+1\).
\end{proof}

\begin{remark}
Theorem~\ref{thm:first} establishes a lower bound.
Exact values of \(j(Q_n)\) and minimal orders of quandles
with prescribed indices are not determined here.
\end{remark}

\section{Second family: unbounded reductivity}\label{sec:second}

\subsection{The reduced free group and a finite quotient}

Let \(F(x_1,\ldots,x_m)\) be the free group.
Define the reduced free group
\[
 RF_m=F(x_1,\ldots,x_m)/
 \langle\!\langle [x_r,x_r^w]:
 1\le r\le m,\ w\in F(x_1,\ldots,x_m)\rangle\!\rangle.
\]
Each \(x_r\) commutes with all its conjugates.
By \cite[Propositions~1.3 and~1.18]{Darne},
\[
 \gamma_{m+1}(RF_m)=1,\qquad
 \gamma_m(RF_m)=Z(RF_m)\cong\Z^{(m-1)!}.
\]
Here and below, the numbering in~\cite{Darne} follows arXiv version v2.
Thus \(RF_m\) has exact class \(m\).
Choose \(1\ne z_m\in\gamma_m(RF_m)\).

We justify the existence of a finite quotient in which \(z_m\) survives.
Consider the ring
\[
 \mathscr A_m=\Z\langle X_1,\ldots,X_m\rangle/I,
\]
where \(I\) is generated by words with at least one repeated letter.
Its additive basis consists of the empty word and all words without repetitions.
For the ideal \(J=(X_1,\ldots,X_m)\), we have \(J^{m+1}=0\);
thus \(1+J\) is a group, with inverses
\[
 (1+u)^{-1}=\sum_{t=0}^{m}(-u)^t.
\]
By \cite[Corollary~1.13]{Darne}, the Magnus map
\[
 \mu:RF_m\hookrightarrow1+J,\qquad x_r\mapsto1+X_r
\]
is injective.
The expansion \(\mu(z_m)-1=\sum_w c_ww\)
has a nonzero integer coefficient \(c_w\).
Choose a prime \(p\) not dividing this coefficient
and reduce the ring modulo \(p\).
The group \(1+\overline J\) is finite, since \(\overline J\)
has a finite basis over \(\F_p\); its order is \(p^{d_m}\), where
\[
 d_m=\sum_{t=1}^{m}\frac{m!}{(m-t)!}.
\]
Let \(A_m\) be the image of \(RF_m\) in this group,
and let \(\phi_m:RF_m\twoheadrightarrow A_m\) be the epimorphism onto the image.
Then \(\phi_m(z_m)\ne1\).
Epimorphisms preserve the lower central series:
\(\phi_m(\gamma_t(RF_m))=\gamma_t(A_m)\),
as follows by induction from preservation of commutators.
Hence \(A_m\) has exact class \(m\).

To make the conjugacy classes of the chosen generators distinct, set
\[
 \alpha:RF_m\to(\Z/2\Z)^m,\quad x_r\mapsto e_r,\qquad
 G_m=\operatorname{im}(\phi_m,\alpha)
 \le A_m\times(\Z/2\Z)^m.
\]
The map \(\alpha\) is well-defined because all the relations
lie in the commutator subgroup of the free group.
Denote the epimorphism onto \(G_m\) by \(\psi_m\),
and the images of \(x_r\) by \(a_r\).
The group \(G_m\) is finite, is generated by \(a_1,\ldots,a_m\),
and has exact class \(m\):
it is a quotient of \(RF_m\), and its first coordinate preserves \(z_m\).
The abelian coordinates \(e_r\) are distinct,
so \(a_r,a_s\) cannot be conjugate for \(r\ne s\).

\subsection{Orbits and the three exact indices}

\begin{proof}[Proof of Theorem~\ref{thm:second}]
Fix \(m\ge2\), and write \(G=G_m\) for brevity.
Set
\[
 C_r=a_r^G,\qquad P_m=\bigcup_{r=1}^{m}C_r,\qquad
 u\conj v=u^{-1}vu.
\]
By Lemma~\ref{lem:orbits}, this is a finite quandle.
The classes \(C_r\) are pairwise disjoint.
The relations of the reduced free group give
\([a_r,a_r^g]=1\) for every \(g\in G\).
Any two elements \(a_r^g,a_r^h\) commute:
conjugation by \(g^{-1}\) sends them to \(a_r,a_r^{hg^{-1}}\).
Thus every \(C_r\) is trivial as a quandle.

Consider the homomorphism
\[
 \rho:G\to\operatorname{Sym}(P_m),\qquad
 \rho(g)(v)=gvg^{-1}.
\]
With the usual composition of permutations, this is a homomorphism;
the maps \(v\mapsto g^{-1}vg\) form an antihomomorphism.
We have \(L_u=\rho(u^{-1})\).
Since \(P_m\) contains all the \(a_r\), it generates \(G\);
hence \(\Inn(P_m)=\rho(G)\).
Its orbits are exactly \(C_1,\ldots,C_m\).
Every orbit series reaches one of the \(C_r\) at its first step
and a singleton quandle at its second.
Thus \(P_m\in\tOS_2\).
The quandle is not trivial: otherwise all the \(a_r\) would
commute pairwise, and the group \(G\) that they generate would be abelian,
contrary to its class \(m\ge2\).
Therefore \(j(P_m)=2\).

The element \(u\conj v\) lies in the same \(C_r\) as \(v\),
and therefore commutes with \(v\).
It follows that \((u\conj v)\conj v=v\), so \(P_m\in\LR_2\).
Nontriviality gives \(k(P_m)=2\).

The kernel of \(\rho\) is \(C_G(P_m)=Z(G)\),
since \(P_m\) generates \(G\).
Consequently,
\[
 \Inn(P_m)\cong G/Z(G).
\]
The quotient \(G/Z(G)\) has exact class \(m-1\).
Indeed, \(\gamma_m(G)\le Z(G)\), so \(\gamma_m(G/Z(G))=1\).
If \(\gamma_{m-1}(G/Z(G))=1\), then
\(\gamma_{m-1}(G)\le Z(G)\), forcing \(\gamma_m(G)=1\),
a contradiction.
Proposition~\ref{prop:criterion} now gives \(i(P_m)=m\).
All the required equalities follow.
\end{proof}

\begin{remark}
The abelian factor \((\Z/2\Z)^m\) is needed to ensure
exactly \(m\) distinct orbits.
For unboundedness of \(i\) with \(j=2\), \(A_m\) itself is sufficient.
The prime \(p\) is chosen after choosing \(z_m\);
we do not assert that an arbitrary \(z_m\) survives
for any prime \(p\) fixed in advance.
\end{remark}

\subsection{\texorpdfstring{An example for \(m=2\)}{An example for m=2}}

To illustrate the index shift, one may take \(G=UT_3(\F_2)\)
with generators
\[
 a=1+E_{12},\qquad b=1+E_{23},\qquad z=1+E_{13}.
\]
This group has order \(8\), class \(2\), and center \(\{1,z\}\).
The generator classes are \(\{a,az\}\) and \(\{b,bz\}\).
On \(P=\{a,az,b,bz\}\), the operation has table
\[
 \begin{array}{c|cccc}
 \conj&a&az&b&bz\\ \hline
 a&a&az&bz&b\\
 az&a&az&bz&b\\
 b&az&a&b&bz\\
 bz&az&a&b&bz
 \end{array}
\]
There are two trivial orbits of size \(2\),
and \(\Inn(P)\cong G/Z(G)\) is nontrivial and abelian.
Thus \((k(P),j(P),i(P))=(2,2,2)\).

\section{Scope of the results}

Both obstructions occur already at the least nontrivial
value of the controlling index, namely \(2\).
The first family consists of a single group conjugacy class;
the second consists of a union of the classes of chosen generators.
It is not the full conjugation quandle on the entire group.
Consequently, \cite[Corollary~5.24]{BCNW}, which concerns
the full quandle \(\operatorname{Conj}(G)\),
does not bound the class of the group in the construction of \(P_m\).

The finite quotient in the second family is chosen
to preserve a nonzero commutator of weight \(m\).
An arbitrary finite quotient may decrease the class
and does not ensure the required index.
The proofs establish the existence of these families;
optimal orders and the exact depth of the first family
are outside the scope of this paper.

\begin{thebibliography}{9}
\bibitem{BCNW}
M.~Bonatto, A.~S.~Crans, T.~Nasybullov, G.~T.~Whitney.
\newblock Quandles with orbit series conditions.
\newblock \emph{Journal of Algebra}, \textbf{567} (2021), 284--309.
\newblock DOI:
\href{https://doi.org/10.1016/j.jalgebra.2020.09.026}{10.1016/j.jalgebra.2020.09.026}.
\newblock Numbering:
\href{https://arxiv.org/abs/2004.10227v3}{arXiv:2004.10227v3}.

\bibitem{NTT}
M.~Noce, G.~Tracey, G.~Traustason.
\newblock A left \(3\)-Engel element whose normal closure is not nilpotent.
\newblock \emph{Journal of Pure and Applied Algebra},
\textbf{224}, no.~3 (2020), 1092--1101.
\newblock DOI:
\href{https://doi.org/10.1016/j.jpaa.2019.07.004}{10.1016/j.jpaa.2019.07.004}.
\newblock \href{https://arxiv.org/abs/1811.12074v1}{arXiv:1811.12074v1}.

\bibitem{Darne}
J.~Darn\'e.
\newblock Milnor invariants of braids and welded braids up to homotopy.
\newblock \emph{Algebraic \& Geometric Topology},
\textbf{24}, no.~3 (2024), 1277--1320.
\newblock DOI:
\href{https://doi.org/10.2140/agt.2024.24.1277}{10.2140/agt.2024.24.1277}.
\newblock Numbering:
\href{https://arxiv.org/abs/1904.10677v2}{arXiv:1904.10677v2}.
\end{thebibliography}
\end{document}
